% Pollis drill-down: Customisation, the settings of Pollis: every setting of Pollis's own and a few others that matter
% most, each with its default, its alternatives and where it is changed (21 November 2026). For users who want to make
% Pollis their own, and for whoever maintains Pollis.
% Compile with: latexmk -pdf workflow_customisation.tex
% Shared look for the one-page Pollis workflow handouts (same fonts and colours as the brochure).
\documentclass[10pt,a4paper]{article}

\usepackage[a4paper,left=17mm,right=17mm,top=14mm,bottom=16mm,footskip=8mm]{geometry}
\usepackage[T1]{fontenc}
\usepackage[default,semibold]{sourcesanspro}
\usepackage{microtype}
\usepackage{graphicx}
\usepackage[table]{xcolor}
\usepackage{tikz}
\usepackage[most]{tcolorbox}
\usepackage{fontawesome5}
\usepackage{enumitem}
\usepackage{fancyhdr}
\usepackage[hidelinks]{hyperref}

\graphicspath{{images/}}

% ---------------------------------------------------------------- colours (as in the brochure)
\definecolor{navy}{HTML}{1C3A5E}
\definecolor{pblue}{HTML}{3B78B5}
\definecolor{amber}{HTML}{F2B33D}
\definecolor{deepamber}{HTML}{D9902A}
\definecolor{pale}{HTML}{EEF3F9}
\definecolor{ink}{HTML}{2A2F36}
\definecolor{logonavy}{HTML}{173D63}
\definecolor{logodark}{HTML}{C98A1E}
\definecolor{logoamber}{HTML}{F2B446}
\definecolor{logocream}{HTML}{FEEFD2}
\definecolor{soft}{HTML}{5E6773}
\color{ink}

\setlength{\parindent}{0pt}
\setlength{\parskip}{3pt}
\setlist{nosep,leftmargin=*,itemsep=2pt}

\pagestyle{fancy}
\fancyhf{}
\renewcommand{\headrulewidth}{0pt}
\fancyfoot[L]{\footnotesize\color{soft}\textbf{Pollis}\enspace Bright Insights through Exploration, Modelling, Simulation and Optimisation}
\fancyfoot[R]{\footnotesize\color{soft}github.com/saragga/pollis/releases/latest}

% The Pollis logo, redrawn so it prints sharply at any size
\newcommand{\pollislogo}[1]{%
  \begin{tikzpicture}[x=#1/512,y=-#1/512,baseline=(current bounding box.center)]
    \fill[logonavy] (0,0) rectangle (512,512);
    \fill[logodark] (36,422) -- (104,256) -- (172,422) -- cycle;
    \fill[logoamber] (104,422) -- (203,180) -- (302,422) -- cycle;
    \fill[logocream] (203,422) -- (339,89) -- (476,422) -- cycle;
  \end{tikzpicture}}

% Top of the page: logo, name, kind of document (Workflow unless given, e.g. [Drill-Down]) and its subject;
% then the question and a one-line promise
\newcommand{\workflowhead}[4][Workflow]{%
  \thispagestyle{fancy}%
  \begin{minipage}[c]{13mm}\pollislogo{11mm}\end{minipage}%
  \begin{minipage}[c]{\dimexpr\linewidth-13mm\relax}
    {\color{navy}\Large\bfseries Pollis}\enspace{\color{soft}#1}\hfill
    {\color{deepamber}\small\bfseries\MakeUppercase{#2}}
  \end{minipage}\par\vspace{7mm}
  {\color{navy}\fontsize{22}{26}\selectfont\bfseries #3}\par\vspace{4pt}
  {\color{soft}\large #4}\par\vspace{5mm}}

% A menu path: \menu{Explore}{Access Data Libraries}... is written as \path{Explore \gt Access Data Libraries}
\newcommand{\gt}{\,{\color{deepamber}\bfseries ›}\,}
\newcommand{\menupath}[1]{{\footnotesize\bfseries\color{navy}#1}}

% A screenshot if images/<file> exists, otherwise an orange placeholder that says what to capture
\newlength{\shotheight}\setlength{\shotheight}{43mm}
\newcommand{\stepshot}[2]{%
  \IfFileExists{images/#1}{%
    \tcbox[enhanced, colback=black, colframe=black!70, boxrule=0.4pt, arc=2pt,
      left=0pt,right=0pt,top=0pt,bottom=0pt, boxsep=0pt, halign=center]{%
      \includegraphics[width=\dimexpr\linewidth-1pt\relax,height=\shotheight,keepaspectratio]{#1}}}{%
    \begin{tcolorbox}[enhanced, colback=white, colframe=orange!80, boxrule=0.6pt, arc=2pt,
      height=\shotheight, valign=center, halign=center, left=6pt, right=6pt]
      {\color{orange}\Large\faCamera}\par\smallskip
      {\color{orange}\footnotesize\bfseries Screenshot: #2}\par\smallskip
      {\color{orange!80!black}\scriptsize images/#1}
    \end{tcolorbox}}}

% One step: number, title, menu path, image file, what to capture, what it means
\newcommand{\step}[6]{%
  \begin{tcolorbox}
    {\color{deepamber}\Large\bfseries #1}\enspace{\color{navy}\large\bfseries #2}\par\vspace{1pt}
    \menupath{#3}\par\vspace{4pt}
    \stepshot{#4}{#5}\par\vspace{3pt}
    {\small #6}
  \end{tcolorbox}}

\newenvironment{steps}{%
  \begin{tcbraster}[raster columns=2, raster equal height=rows, raster column skip=4mm,
    raster row skip=4mm, enhanced, colback=pale, colframe=pale, boxrule=0pt, sharp corners,
    left=6pt,right=6pt,top=6pt,bottom=6pt, fontupper=\raggedright]}{\end{tcbraster}}

% Boxes for the bottom of the page
\newtcolorbox{pbox}[2][]{enhanced, colback=pale, colframe=amber, boxrule=0pt,
  leftrule=3pt, sharp corners, left=7pt, right=7pt, top=6pt, bottom=6pt,
  fontupper=\small\raggedright, before upper={{\color{navy}\normalsize\bfseries #2}\par\smallskip}, #1}
\newtcolorbox{nbox}[1][]{enhanced, colback=navy, colframe=navy, boxrule=0pt,
  sharp corners, left=8pt, right=8pt, top=6pt, bottom=6pt, coltext=white,
  fontupper=\small\raggedright, #1}

% A Next Steps item inside the bottom boxes
\newcommand{\nsitem}[2]{\item[\color{deepamber}\faChevronRight] \textbf{#1}: #2}

% Citations: \dcite{key}{Breiman (2001)} links to its entry \dref{key}{...} in the References, keyed as in pollis.bib
\newcommand{\dcite}[2]{\hyperlink{ref:#1}{\color{pblue}#2}}
\newcommand{\dref}[2]{\raisebox{1.1em}[0pt][0pt]{\hypertarget{ref:#1}{}}#2\par}
\newcommand{\ddoi}[1]{\href{https://doi.org/#1}{\color{pblue}doi:#1}}
\newcommand{\references}{\section*{References}\addcontentsline{toc}{section}{References}\footnotesize\raggedright\setlength{\parskip}{3pt}}

% The page of each heading, for Ask Pollis
% Pollis: the page of each heading, for Ask Pollis. A document that inputs this file, at the end of its preamble, writes
% <jobname>.labels beside its PDF, a line page|level|heading for each heading, the page counted from the first, the
% cover, as 1. compile_documents.jl copies it to ~/.pollis/documents, and ask_index.jl cuts the document into passages
% under its headings with it. \asklabel{1}{Title} marks a topic or a page, \asklabel{2}{Title} a heading inside it; the
% sections and subsections of a document that has them are marked as they go into its contents. A package vignette
% inputs it before its \documentclass (compile_vignette.jl), so the contents are taken over when the document begins,
% after the packages have taken them over. Read a second time, as by a copy of a document for Edit the LaTeX Source,
% which inputs it as resources/ask_labels when Pollis has put it in already (document_source.jl), it does nothing.
\ifdefined\asklabel\endinput\fi
\makeatletter
\newwrite\pollis@labels
\immediate\openout\pollis@labels=\jobname.labels
\newcommand{\asklabel}[2]{\write\pollis@labels{\the\ReadonlyShipoutCounter|#1|\unexpanded{#2}}}
\def\pollis@toc{toc}
\newcommand{\pollis@contentsline}[3]{%
  \def\pollis@file{#1}%
  \ifx\pollis@file\pollis@toc\asklabel{#2}{#3}\fi
  \pollis@addcontentsline{#1}{#2}{#3}}
% The headings of the New Methods style of the vignettes, \nmsection and \nmsub, which go into no contents
\def\pollis@asksection#1{\pollis@nmsection{#1}\asklabel{1}{#1}}
\def\pollis@asksub#1{\pollis@nmsub{#1}\asklabel{2}{#1}}
\AddToHook{begindocument/end}{%
  \let\pollis@addcontentsline\addcontentsline\let\addcontentsline\pollis@contentsline
  \ifdefined\nmsection\let\pollis@nmsection\nmsection\let\nmsection\pollis@asksection\fi
  \ifdefined\nmsub\let\pollis@nmsub\nmsub\let\nmsub\pollis@asksub\fi}
\makeatother



\usepackage{titlesec}
\usepackage{tabularx}
\usepackage{booktabs}
\usepackage{longtable}

\fancyhf{}
\fancyfoot[L]{\footnotesize\color{soft}\textbf{Pollis}\enspace Drill-down: Customisation}
\fancyfoot[R]{\footnotesize\color{soft}\thepage}

\titleformat{\section}{\color{navy}\large\bfseries}{\color{deepamber}\thesection}{0.6em}{}
\titleformat{\subsection}{\color{navy}\normalsize\bfseries}{\color{deepamber}\thesubsection}{0.5em}{}
\titlespacing*{\section}{0pt}{12pt}{4pt}
\titlespacing*{\subsection}{0pt}{8pt}{2pt}
\setlength{\parskip}{4pt}

\usepackage{xurl}
\newcommand{\file}[1]{{\small\nolinkurl{#1}}}
\newcolumntype{L}{>{\raggedright\arraybackslash}X}
\newcommand{\ui}[1]{\textbf{#1}}
% Code with braces or spaces in it, which \file would not show as written
\newcommand{\code}[1]{\texttt{\small #1}}
% A setting: its name in the Settings editor, its id and its default
\newcommand{\setitem}[3]{\item \ui{#1} (\file{#2}), \emph{#3}:}
% One theme's picture with its name
\newcommand{\themeshot}[2]{%
  \begin{minipage}[t]{0.82\linewidth}\centering
    \includegraphics[width=\linewidth]{#1}\\[2pt]
    {\small\ui{#2}}
  \end{minipage}}

\begin{document}

\workflowhead[Drill-Down]{Customisation}
  {The settings of Pollis}
  {Every setting of Pollis's own, and a few others that matter most: what each does, its default, the alternatives,
   and where to change it.}

\begin{nbox}
{\color{amber}\bfseries Drill-down.}\enspace
Pollis works as it is installed. Its settings change how it looks, where charts and maps appear, which models answer
questions and run agents, and which feeds it reads. Most of them have a page or a command of their own in Pollis, and
all of them can be changed in the Settings editor.

This drill-down lists every setting of Pollis's own, also those that are empty until you use the feature, and a few
settings of the editor and of Julia that matter most. It describes the state of 21 November 2026; \S\ref{sec:glance}
lists them all at a glance.
\end{nbox}

\tableofcontents

\section{Changing a setting}\label{sec:change}
\begin{itemize}
  \item \textbf{The Settings editor.} Choose the \ui{Manage} gear at the bottom left, then \ui{Settings}, or press
    \ui{Cmd+,} (\ui{Ctrl+,} on Windows and Linux). Type in the search box a word, or a setting's id, such as
    \file{pollis.plots.keepLatest}. A setting changed from its default has a bar beside it, and its gear menu has
    \ui{Reset Setting}.
  \item \textbf{The settings file.} \ui{Preferences: Open User Settings (JSON)}, in the Command Palette, opens
    \file{settings.json}, where each setting is its id and its value. Settings that are lists, such as the RSS feeds,
    are edited there.
  \item \textbf{User and Workspace.} A setting of the \ui{User} tab applies to every window; one of the
    \ui{Workspace} tab applies only to the folder open. Most settings of Pollis apply to the whole application, and
    are in the \ui{User} tab only.
  \item \textbf{In Pollis itself.} Most settings have a page, a command or a button that changes them, named below
    with each. A change there is written to the setting, and a change of the setting shows there.
\end{itemize}

\section{Appearance}\label{sec:appearance}
\subsection{Colour theme}
\begin{itemize}
  \setitem{Workbench: Color Theme}{workbench.colorTheme}{Pollis 2026 Dark}
    the colours of the whole window and of the code. Change it with \ui{Manage > Themes > Color Theme}, which shows
    each theme as you move through the list. Pollis has four themes of its own, at the top of the list: the dark ones
    above the other dark themes, the light ones above the other light themes.
\end{itemize}
\begin{center}
\themeshot{images/theme_pollis_dark.png}{Pollis Dark}\\[8pt]
\themeshot{images/theme_pollis_2026_dark.png}{Pollis 2026 Dark}
\end{center}
\begin{itemize}
  \item \textbf{Pollis Dark and Pollis Light} are the colours of Dark (Visual Studio) and Light (Visual Studio), with
    Pollis's yellow on four things: the badges with a number, such as those of the \ui{Extensions} view; the line
    beside the selected item of the activity bar; the top of the selected tab, such as \ui{Particle Filter}; and the
    line under the selected title of a panel, such as \ui{Terminal} and \ui{Ask Pollis}.
  \item \textbf{Pollis 2026 Dark and Pollis 2026 Light}, the default themes, have Pollis's own colours for the window,
    darker and with outlined fields, with the same yellow, and the code in the colours of Dark (Visual Studio) and
    Light (Visual Studio).
  \item \textbf{Pollis Light and Pollis 2026 Light} are the light themes of each pair, with the same yellow.
  \item \textbf{The other themes} are those of VS Code, among them the high-contrast themes for low vision, and those
    of extensions you install.
\end{itemize}

\subsection{Pollis Views}
\begin{itemize}
  \setitem{Pollis: View}{pollis.view}{Max}
    the View in use: which menu items Pollis shows. \ui{Max} shows every item; \ui{Min} shows a short selection; your
    own Views are JSON files in \file{~/.pollis/views}. Change it with the name of the View in the title bar, such as
    \ui{View: Max}, or the command \ui{Pollis: Select View...}. The Pollis Guide describes how to make a View.
\end{itemize}

\section{Julia}\label{sec:julia}
Settings of the Julia extension, which runs the Julia REPL, notebooks and the panels' code.
\begin{itemize}
  \setitem{Julia: Num Threads}{julia.NumThreads}{auto}
    how many threads Julia starts with. \file{auto} starts one thread for each core of the computer; a number starts
    that many. It applies when the REPL starts again.
  \setitem{Julia: Start REPL On Startup}{julia.startREPLOnStartup}{on}
    whether the Julia REPL starts with Pollis, so that it is ready when you run a panel's code. Off, Pollis starts
    sooner, and the REPL starts the first time it is needed.
  \setitem{Julia: Executable Path}{julia.executablePath}{empty}
    the Julia that Pollis runs. Empty, it is the \file{julia} of the system's path, which juliaup installs; otherwise
    the path of a Julia, or a channel of juliaup, such as \file{julia +release}.
  \setitem{Julia: Use Revise}{julia.useRevise}{on}
    whether the REPL loads Revise.jl, so that a change to a package or a file you included applies without starting
    the REPL again.
\end{itemize}

\section{Plots and maps}\label{sec:plots}
Where the charts and maps shown in the Julia REPL appear. Notebooks and Pluto notebooks draw them in the cell's output,
whatever these settings say. The drill-down Plotting says what each place costs in memory.
\begin{itemize}
  \setitem{Plots: Show Charts In}{pollis.plots.showChartsIn}{Plots Tab}
    where a chart appears.
    \begin{itemize}
      \item \ui{Plots Tab}: the \ui{Plots} tab, which draws one chart at a time and keeps the others, with the
        \ui{Plot Navigator} and the commands to page through them and save them.
      \item \ui{Separate Window}: the \ui{Plots} tab in a window of its own, which can go on another screen.
      \item \ui{Browser Page}: each chart on a page of its own in Pollis's browser, so that charts can be side by side.
      \item \ui{External Viewer}: each chart on a page of its own in the program of \ui{Plots: External Viewer}.
    \end{itemize}
    A page has none of the commands of the \ui{Plots} tab: a chart is saved from the page.
  \setitem{Plots: Show Maps In}{pollis.plots.showMapsIn}{Browser Page}
    where a map of \file{Pollis.webmap} or \file{Pollis.streetview} appears: \ui{Browser Page}, a page in Pollis's
    browser for each map title, so that maps can be side by side; \ui{Plots Tab}, one map at a time, lighter on
    memory; or \ui{External Viewer}.
  \setitem{Plots: External Viewer}{pollis.plots.externalViewer}{empty}
    the program of \ui{External Viewer}: a web browser, since a chart or a map is a web page. Empty, it is the
    system's default browser; on macOS, otherwise, the name of an application, such as Firefox, Google Chrome or
    Safari, and elsewhere its command or path.
  \setitem{Plots: Keep Latest}{pollis.plots.keepLatest}{50}
    how many plots the \ui{Plots} tab keeps, the oldest removed when a new one would pass it; and how many pages of
    charts Pollis keeps for \ui{Browser Page} and \ui{External Viewer}. Fewer plots use less memory.
\end{itemize}

\section{Ask Pollis}\label{sec:ask}
The settings of \ui{Help > Ask Pollis}. Most are set on its own pages: \ui{Compose > Local AI Models > Set Active
Models} chooses the models, and \ui{Strategy}, under the question, the answering strategy.
\begin{itemize}
  \setitem{Answer With}{pollis.ask.answerWith}{This Computer}
    who writes the answers: \ui{This Computer}, with the chat model of \file{pollis.ask.chatModel}, or a cloud model
    of Anthropic, Google or OpenAI, with your own key, kept in Pollis's secret storage. A cloud model is sent the
    question, the conversation so far and the passages found on this computer, never your files. \ui{Set Models} in
    \ui{Set Active Models} chooses it.
  \setitem{Chat Model}{pollis.ask.chatModel}{the default chat model}
    the model that writes the answers on this computer: a GGUF file of Hugging Face, as
    \file{hf:owner/repository/file.gguf}, downloaded the first time, or the path of a GGUF file. It runs in the
    llama.cpp that comes with Pollis.
  \setitem{Embedding Model}{pollis.ask.embeddingModel}{the default embedding model}
    the model that finds passages by their meaning, in the same forms. When it changes, the passages are embedded
    again the next time the Julia REPL starts; until then they are found by their words only.
  \setitem{Reranker}{pollis.ask.reranker}{empty}
    the model that chooses among the passages found, also for a collection of Retrieval-Augmented Generation without a
    reranker of its own. Empty, there is none.
  \setitem{Strategy}{pollis.ask.strategy}{panelAware}
    how the passages are found and chosen, and how the model is asked: \file{base}, \file{panelAware}, or a strategy
    of yours, a TOML file in \file{~/.pollis/ask/strategies}. \ui{Compose > Local AI Models > Evaluate Answering
    Models...} compares them.
  \setitem{Passages}{pollis.ask.passages}{8}
    how many passages go to the model with each question, from 1 to 20. More passages answer more, but take longer on
    this computer and cost more with a cloud model.
  \setitem{Julia Manual}{pollis.ask.juliaManual}{on}
    whether Ask Pollis searches the Julia Manual too: the Manual, Base and the standard library of the Julia that
    Pollis runs.
  \setitem{DuckDB Manual}{pollis.ask.duckdbManual}{on}
    whether it searches the DuckDB Manual too, at the version of the DuckDB that Pollis carries, downloaded once from
    duckdb.org.
  \setitem{Package APIs}{pollis.ask.packageApis}{on}
    whether it can search the docstrings of the packages of the panels, some twenty thousand passages, about 120 MB in
    the index. They are searched only with \ui{Search In: Package APIs} or \ui{All Sources}.
  \setitem{Passage Prefix}{pollis.ask.passagePrefix}{the form of the default embedding model}
    what each passage is given to the embedding model as: \code{\{title\}} is its document and heading,
    \code{\{text\}} its text. Another embedding model may want another form, or \code{\{text\}} alone.
  \setitem{Query Prefix}{pollis.ask.queryPrefix}{the form of the default embedding model}
    what a question is given to the embedding model as, \code{\{text\}} being the question.
\end{itemize}

\section{Local models}\label{sec:local}
\begin{itemize}
  \setitem{Local Models: Publishers}{pollis.localModels.publishers}{ten publishers}
    the publishers of Hugging Face whose GGUF models the \ui{Choose...} button beside a local model field offers:
    small models of the last year, the most downloaded in the last month first. The default list is ggml-org, unsloth,
    bartowski, lmstudio-community, Qwen, google, nomic-ai, mistralai, microsoft and ibm-granite.
\end{itemize}

\section{Agents and models}\label{sec:agents}
The coding agents and the models of the agent tasks: \ui{Compose > Orchestration > Model Tasks and Routing} and
\ui{Compose > Coding Agents} change them all.
\begin{itemize}
  \setitem{Tier Cells}{pollis.agents.tierCells}{Light, Standard and Strong on Claude Code}
    the agent and the model of each tier. Each tier, \file{light}, \file{standard}, \file{strong} or one of your own,
    has its \file{agent}, by name, and its \file{model}: \file{cloud:} and a model of \file{pollis.agents.cloudModels},
    or \file{local:} and a GGUF file served on this computer by the llama.cpp that comes with Pollis.
  \setitem{Cloud Models}{pollis.agents.cloudModels}{the cloud models Pollis lists}
    the cloud models of the rows of Model Tasks and Routing: each its \file{name} and, for each agent that runs it,
    the arguments that choose it, such as \code{"Codex": "-m model-name"}. An agent with no arguments for a model cannot
    run it.
  \setitem{Task Tiers}{pollis.agents.taskTiers}{empty}
    the tier of an agent task when it is not its default one.
  \setitem{Other Agents}{pollis.agents.otherAgents}{empty}
    agents of your own, besides Claude Code, Codex, OpenCode, GitHub Copilot CLI and Antigravity CLI: each its name
    and its command line, with \code{\{prompt\}} where the first message goes.
\end{itemize}

\section{Literature}\label{sec:literature}
The feeds of \ui{Explore > RSS > RSS Feeds} and of the welcome page. All are set from \ui{RSS Feeds} itself; the
lists are kept in these settings.
\begin{itemize}
  \setitem{Feeds}{pollis.rss.feeds}{empty}
    the RSS and Atom feeds followed: an address, or an address with a topic and a name. \ui{Add Feed} adds one.
  \setitem{Display Days}{pollis.rss.displayDays}{30}
    how far back items are shown: the last 7, 30 or 90 days, or 0 for every item kept. Items without a date are
    always shown.
  \setitem{Alerts Mailbox}{pollis.rss.alertsMailbox}{empty}
    the mailbox the email alerts of Google Scholar, JSTOR, HeinOnline and other sources are sent to: its IMAP server,
    port (993), user and folder (\file{INBOX}). Pollis reads it, read only, to make a feed of each alert. Set it with
    \ui{Add Feed > Set Mailbox...}, which keeps the password in the system's keychain.
  \setitem{Hidden Feeds}{pollis.rss.hiddenFeeds}{empty}
    the feeds whose items are not shown, in RSS Feeds or on the welcome page; they are still read and kept.
  \setitem{Welcome Hidden Feeds}{pollis.rss.welcomeHiddenFeeds}{empty}
    the feeds whose news the welcome page does not show.
  \setitem{Unsubscribed Feeds}{pollis.rss.unsubscribedFeeds}{empty}
    the feeds no longer read; their items stay, and are still shown.
  \setitem{Removed Feeds}{pollis.rss.removedFeeds}{empty}
    the feeds removed from RSS Feeds, which \ui{Restore} under \ui{Add Feed > Removed Feeds} brings back.
\end{itemize}

\section{Workspace, content and publishing}\label{sec:workspace}
\begin{itemize}
  \setitem{Pollis Workspace: Location}{pollis.workspace.location}{empty}
    the folder of your Pollis Workspace, where Pollis saves your work and the copies you choose to save, also when the
    workspace is not open in the window. Pollis sets it when you create or open a Pollis Workspace; one open in the
    window is used first.
  \setitem{Content: Check for Updates}{pollis.content.checkForUpdates}{on}
    whether Pollis checks once a day for new and updated replications and published benchmarks, published between
    releases of Pollis.
  \setitem{Content: Repository}{pollis.content.repository}{\file{https://github.com/Trumpingtons/pollis-content}}
    the GitHub repository that content is read from.
  \setitem{Article Methods: GitHub Owner}{pollis.newMethods.githubOwner}{Trumpingtons}
    the GitHub user or organisation that \ui{Package and Publish}, on the page of an article method, pushes method
    packages to.
  \setitem{Article Methods: Registry}{pollis.newMethods.registry}{\file{https://github.com/Trumpingtons/PollisRegistry}}
    the Julia registry that \ui{Package and Publish} registers method packages in, when they meet the rules of the
    Pollis registry.
  \setitem{Collaboration: Signalling}{pollis.collaboration.signalling}{nostr}
    the free public servers that introduce the participants of a session you start, after which they connect
    directly: \file{nostr}, public nostr relays, or \file{mqtt}, public MQTT brokers, for networks that block the
    relays. The invitation says which, so the others need not change it. \ui{Compose > Collaboration > Test
    Connection...} shows which are reachable from your network.
\end{itemize}

\section{All settings at a glance}\label{sec:glance}
{\small
\renewcommand{\arraystretch}{1.15}
\begin{longtable}{@{}p{0.43\linewidth}p{0.2\linewidth}p{0.31\linewidth}@{}}
\toprule
\textbf{Id} & \textbf{Default} & \textbf{Changed in Pollis with}\\
\midrule
\endhead
\multicolumn{3}{@{}l}{\textbf{Appearance}}\\
\file{workbench.colorTheme} & Pollis 2026 Dark & Manage > Themes > Color Theme\\
\file{pollis.view} & Max & Pollis: Select View...\\
\multicolumn{3}{@{}l}{\textbf{Julia}}\\
\file{julia.NumThreads} & auto & Settings\\
\file{julia.startREPLOnStartup} & on & Settings\\
\file{julia.executablePath} & empty & Settings\\
\file{julia.useRevise} & on & Settings\\
\multicolumn{3}{@{}l}{\textbf{Plots and maps}}\\
\file{pollis.plots.showChartsIn} & Plots Tab & Settings\\
\file{pollis.plots.showMapsIn} & Browser Page & Settings\\
\file{pollis.plots.externalViewer} & empty & Settings\\
\file{pollis.plots.keepLatest} & 50 & Settings\\
\multicolumn{3}{@{}l}{\textbf{Ask Pollis}}\\
\file{pollis.ask.answerWith} & This Computer & Set Active Models\\
\file{pollis.ask.chatModel} & default chat model & Set Active Models\\
\file{pollis.ask.embeddingModel} & default embedding model & Set Active Models\\
\file{pollis.ask.reranker} & empty & Set Active Models\\
\file{pollis.ask.strategy} & panelAware & Strategy, under the question\\
\file{pollis.ask.passages} & 8 & Settings\\
\file{pollis.ask.juliaManual} & on & Settings\\
\file{pollis.ask.duckdbManual} & on & Settings\\
\file{pollis.ask.packageApis} & on & Settings\\
\file{pollis.ask.passagePrefix} & default form & Settings\\
\file{pollis.ask.queryPrefix} & default form & Settings\\
\multicolumn{3}{@{}l}{\textbf{Local models}}\\
\file{pollis.localModels.publishers} & ten publishers & Settings (JSON)\\
\multicolumn{3}{@{}l}{\textbf{Agents and models}}\\
\file{pollis.agents.tierCells} & Claude Code & Model Tasks and Routing\\
\file{pollis.agents.cloudModels} & Pollis's list & Model Tasks and Routing\\
\file{pollis.agents.taskTiers} & empty & Model Tasks and Routing\\
\file{pollis.agents.otherAgents} & empty & Coding Agents\\
\multicolumn{3}{@{}l}{\textbf{Literature}}\\
\file{pollis.rss.feeds} & empty & RSS Feeds > Add Feed\\
\file{pollis.rss.displayDays} & 30 & Settings\\
\file{pollis.rss.alertsMailbox} & empty & Add Feed > Set Mailbox...\\
\file{pollis.rss.hiddenFeeds} & empty & RSS Feeds\\
\file{pollis.rss.welcomeHiddenFeeds} & empty & RSS Feeds\\
\file{pollis.rss.unsubscribedFeeds} & empty & RSS Feeds\\
\file{pollis.rss.removedFeeds} & empty & Add Feed > Removed Feeds\\
\multicolumn{3}{@{}l}{\textbf{Workspace, content and publishing}}\\
\file{pollis.workspace.location} & empty & creating or opening a Pollis Workspace\\
\file{pollis.content.checkForUpdates} & on & Settings\\
\file{pollis.content.repository} & pollis-content & Settings\\
\file{pollis.newMethods.githubOwner} & Trumpingtons & Settings\\
\file{pollis.newMethods.registry} & PollisRegistry & Settings\\
\file{pollis.collaboration.signalling} & nostr & Settings\\
\bottomrule
\end{longtable}
}

\section{Where things are}
\begin{center}\small
\renewcommand{\arraystretch}{1.2}
\begin{tabularx}{\linewidth}{@{}p{0.55\linewidth}L@{}}
\toprule
\textbf{File} & \textbf{Settings}\\
\midrule
\file{extensions/theme-defaults/themes/pollis-*.json} & The four Pollis themes\\
\file{src/vs/workbench/contrib/pollisViews/browser/pollisViews.contribution.ts} & \file{pollis.view}\\
\file{extensions/language-julia/package.json} & \file{julia.*}, \file{pollis.plots.*},
  \file{pollis.localModels.publishers}\\
\file{src/vs/workbench/contrib/compose/browser/askPollisPage.ts} & \file{pollis.ask.*}\\
\file{src/vs/workbench/contrib/compose/browser/modelRouting.ts} & \file{pollis.agents.tierCells}, \file{cloudModels},
  \file{taskTiers}\\
\file{src/vs/workbench/contrib/compose/browser/pollisAgents.ts} & \file{pollis.agents.otherAgents}\\
\file{src/vs/workbench/contrib/monitor/browser/rssFeedsService.ts} & \file{pollis.rss.*}\\
\file{src/vs/workbench/contrib/pollisWorkspace/browser/pollisWorkspace.contribution.ts} &
  \file{pollis.workspace.location}\\
\file{src/vs/workbench/contrib/compose/browser/pollisContent.ts} & \file{pollis.content.*}\\
\file{src/vs/workbench/contrib/compose/browser/newMethods.ts} & \file{pollis.newMethods.*}\\
\file{src/vs/workbench/contrib/compose/browser/collaboration/collaboration.contribution.ts} &
  \file{pollis.collaboration.signalling}\\
\bottomrule
\end{tabularx}
\end{center}

\end{document}
